% !TeX program = xelatex
% Póster A0 vertical: 841 x 1189 mm.
% Compile con XeLaTeX (también admite LuaLaTeX).
\documentclass[12pt]{article}
\usepackage[spanish]{babel}
\bibliographystyle{abbrvnat} % Activar al añadir una bibliografía.
\usepackage[numbers, square]{natbib}
\usepackage[
  paperwidth=841mm,paperheight=1189mm,
  left=25mm,right=25mm,top=190mm,bottom=25mm,
  headheight=0pt,headsep=0pt,footskip=0pt
]{geometry}
\usepackage{investigathon}
\usepackage{amsmath}
\usepackage{unicode-math}
\setmathfont{latinmodern-math.otf}
\usepackage{caption}
\DeclareCaptionFont{poster}{\fontsize{24}{30}\selectfont}
\captionsetup{font=poster,labelfont=bf}

% Ajustar esta definición si el conjunto usa otra notación.
\newcommand{\Rs}{\mathcal{R}_{s}}

% Las imágenes se cargan automáticamente cuando están disponibles.
\newcommand{\PosterImage}[2][0.7\linewidth]{%
  \IfFileExists{#2}{%
    \includegraphics[width=#1,keepaspectratio]{#2}%
  }{%
    \begin{tcolorbox}[
      width=#1,height=65mm,
      colback=PosterFill,colframe=PosterHeading!35,
      boxrule=0.6pt,arc=0pt,
      halign=center,valign=center,
      fontupper=\fontsize{22}{28}\selectfont,
      colupper=PosterText
    ]
      Imagen pendiente\\[3mm]
      \texttt{\detokenize{#2}}
    \end{tcolorbox}%
  }%
}
\usepackage{multicol}
\usepackage{tcolorbox}
\usepackage{tikz}
\usetikzlibrary{positioning,babel}
\tcbuselibrary{skins}

% Two equal body columns with a 20 mm gap.
\setlength{\columnsep}{20mm}
\setlength{\columnseprule}{0pt}
\setlength{\multicolsep}{0pt}
\setlength{\emergencystretch}{2em}
\setlength{\parskip}{0.35\baselineskip}

% ==========================================================
% BODY STYLES — local to this document
% Colours follow the dark green / sage tones of the header.
% ==========================================================
\definecolor{PosterHeading}{HTML}{1B3430} % Deep teal
\definecolor{PosterFill}{HTML}{EAF0E9}    % Pale sage
\definecolor{PosterText}{HTML}{1A2827}    % Dark charcoal

\newcommand{\sectionfont}{%
  \normalfont\fontsize{44}{54}\selectfont\bfseries
}
\newcommand{\bodyfont}{%
  \normalfont\fontsize{30}{40}\selectfont
}

% Section headings (visual, unnumbered).
\renewcommand{\section}[1]{%
  \begin{tcolorbox}[
    enhanced,
    colback=PosterHeading,
    colframe=PosterHeading,
    colupper=white,
    fontupper=\sectionfont,
    boxrule=0pt,
    arc=0pt,
    outer arc=0pt,
    left=25pt, right=25pt,
    top=20pt, bottom=20pt,
    drop shadow=black!15,
    before skip=1cm,
    after skip=1cm
  ]
    #1
  \end{tcolorbox}%
}

% Filled poster box: deep teal title and pale sage body.
\newtcolorbox{filledbox}[2][]{
  enhanced,
  before skip=1cm,
  after skip=1cm,
  colback=PosterFill,
  colframe=PosterHeading,
  colbacktitle=PosterHeading,
  coltitle=white,
  colupper=PosterText,
  fonttitle=\sectionfont,
  fontupper=\bodyfont,
  title=#2,
  boxrule=0pt,
  leftrule=0pt, rightrule=0pt, bottomrule=0pt,
  arc=0pt,
  outer arc=0pt,
  left=25pt, right=25pt, top=25pt, bottom=25pt,
  toptitle=20pt, bottomtitle=20pt,
  titlerule=0pt,
  drop shadow=black!15,
  #1
}

% White poster box; retains the border-free geometry supplied.
\newtcolorbox{borderbox}[2][]{
  enhanced,
  before skip=1cm,
  after skip=1cm,
  colback=white,
  colframe=PosterHeading,
  colbacktitle=PosterHeading,
  coltitle=white,
  colupper=PosterText,
  fonttitle=\sectionfont,
  fontupper=\bodyfont,
  title=#2,
  boxrule=0pt,
  leftrule=0pt, rightrule=0pt, bottomrule=0pt,
  arc=0pt,
  outer arc=0pt,
  left=25pt, right=25pt, top=25pt, bottom=25pt,
  toptitle=20pt, bottomtitle=20pt,
  titlerule=0pt,
  drop shadow=black!15,
  #1
}

% Unfilled poster box: white body and deep teal title.
\newtcolorbox{unfilledbox}[2][]{
  enhanced,
  before skip=1cm,
  after skip=1cm,
  colback=white,
  colframe=white,
  colbacktitle=PosterHeading,
  coltitle=white,
  colupper=PosterText,
  fonttitle=\sectionfont,
  fontupper=\bodyfont,
  title=#2,
  boxrule=0pt,
  leftrule=0pt, rightrule=0pt, bottomrule=0pt,
  arc=0pt,
  outer arc=0pt,
  left=25pt, right=25pt, top=25pt, bottom=25pt,
  toptitle=20pt, bottomtitle=20pt,
  titlerule=0pt,
  drop shadow=black!15,
  #1
}

% ==========================================================
% CAMBIE SOLAMENTE ESTOS CUATRO CAMPOS
% track: ia | optimizacion | biologia | fisica-social
% ==========================================================
\PosterSetup{
  titulo={Kuramoto desde la topología},
  integrantes={Tomas Brucellaria, Manuel Ferro, Lucas Grasso Ramos, Tadeo Solari},
  equipo={T08},
  track={optimizacion}
}

\begin{document}
\EncabezadoPoster
\color{PosterText}

% Todo el contenido posterior al encabezado ocupa dos columnas.
\begin{multicols*}{2}
\raggedcolumns
\bodyfont
\setlength{\parskip}{0.35\baselineskip}

% ========================================================
% 1. COLUMNA IZQUIERDA
% ========================================================
\section{¿Por qué topología?}


El problema de energía de Kuramoto \cite{10.1007/BFb0013365} se puede entender viendo los ciclos del grafo.

El camino es un grafo que sincroniza globalmente, mientras que si le agregamos la arista para formar el ciclo, este no sincroniza

\begin{center}
    \begin{tikzpicture}[
        vertice/.style={
            circle,
            draw,
            minimum size=7mm,
            inner sep=0pt,
            line width=0.8pt
        },
        arista/.style={line width=0.8pt}
    ]
    
        % Grafo izquierdo: camino P5
        \foreach \i in {0,...,4} {
            \node[vertice] (p\i) at ({1.2*\i},0) {};
        }
    
        % Grafo derecho: ciclo C5
        \foreach \i in {0,...,4} {
            \node[vertice] (c\i) at ({6.5+1.2*\i},0) {};
        }
    
        % Aristas horizontales
        \foreach \a/\b in {0/1,1/2,2/3,3/4} {
            \draw[arista] (p\a) -- (p\b);
            \draw[arista] (c\a) -- (c\b);
        }
    
        % Arista curva del segundo grafo
        \draw[arista]
            (c0.north) to[out=70,in=110,looseness=1.1]
            (c4.north);
    
    \end{tikzpicture}
\end{center}

% ========================================================
% 2. CONTENIDO EN COLUMNAS
% ========================================================

\begin{filledbox}{Conceptos}
    \begin{center}
        ¿Cuanto se enrolla un ciclo \(\mathcal{C}\)?
        \PosterImage[0.65\linewidth]{assets/reloje.png}
    \end{center}
    \[
        wrap(\theta) = (\theta + \pi) \text{ mod } 2\pi - \pi \in [\pi,\pi)
    \]
    \[
    q(\mathcal{C})=\frac1{2\pi}\sum_{v_i \in \mathcal{C}}wrap(\varphi_{v_{i+1}} - \varphi_{v_i})\in \mathbb{Z} .
    \]
    
    \begin{center}
        \begin{tikzpicture}[scale=1.7,
            % Definición de estilos para mantener el código limpio
            dot/.style={circle, fill=PosterHeading, inner sep=0pt, minimum size=6pt},
            axis/.style={dashed, thin},
            basecirc/.style={thick, color=PosterHeading},
            redarc/.style={line width=2.5pt, color=red!80!black}
        ]
        
        % ================= Primer círculo (Izquierda) =================
        \begin{scope}[xshift=0cm]
            % Ejes
            \draw[axis] (-1.8,0) -- (1.8,0);
            \draw[axis] (0,-1.8) -- (0,1.8);
            % Círculo base
            \draw[basecirc] (0,0) circle (1.5);
            % Arco rojo
            \draw[redarc] (1.5,0) arc (0:45:1.5);
            % Puntos en los extremos
            \node[dot] at (0:1.5) {};
            \node[dot] at (45:1.5) {};
            % Texto inferior
            \node at (0,-2.5) {\fontsize{24}{30}\selectfont $|\cos(\theta_w - \theta_v)| > 0$};
        \end{scope}
        
        % ================= Segundo círculo (Centro) =================
        \begin{scope}[xshift=4.5cm]
            % Ejes
            \draw[axis] (-1.8,0) -- (1.8,0);
            \draw[axis] (0,-1.8) -- (0,1.8);
            % Círculo base
            \draw[basecirc] (0,0) circle (1.5);
            % Arco rojo
            \draw[redarc] (1.5,0) arc (0:140:1.5);
            % Puntos en los extremos
            \node[dot] at (0:1.5) {};
            \node[dot] at (140:1.5) {};
            % Texto inferior
            \node at (0,-2.5) {\fontsize{24}{30}\selectfont $|\cos(\theta_w - \theta_v)| \le 0$};
        \end{scope}
        
        % ================= Tercer círculo (Derecha) =================
        \begin{scope}[xshift=9cm]
            % Ejes
            \draw[axis] (-1.8,0) -- (1.8,0);
            \draw[axis] (0,-1.8) -- (0,1.8);
            % Círculo base
            \draw[basecirc] (0,0) circle (1.5);
            % Arco rojo
            \draw[redarc] (1.5,0) arc (0:170:1.5);
            % Puntos en los extremos
            \node[dot] at (0:1.5) {};
            \node[dot] at (170:1.5) {};
            % Texto inferior (usando align=center para saltos de línea)
            \node[align=center] at (0,-2.8) {\fontsize{24}{30}\selectfont se acerca a \\ \fontsize{24}{30}\selectfont la discontinuidad \\ \fontsize{24}{30}\selectfont de $wrap$};
        \end{scope}
        
        \end{tikzpicture}
    \end{center}
    
    \(\mathcal{R}\): aristas relajadas, a menos de un cuarto de vuelta; forman clusters.
    \(\mathcal{T}\): aristas tensas, a cuarto de vuelta o mas; separan clusters.
    
\end{filledbox}

\begin{borderbox}{Lo que probamos}
    \begin{itemize}
        \item \(q = 0 \quad \forall \mathcal{C} \text{ ciclo } \in G \iff \theta^*\) es un equilibrio sincronizante.
        
        \item Sea $\theta$ un equilibrio no sincronizado de un grafo conexo sin antípodas, con $n\ge3$ y $\mu=\mu(G)$. Luego, si todas las aristas incidentes a $S=\{\phi_1, \phi_2, \phi_3\}$ están en \(\mathcal{R}\), $\mu\le2/3$.
        
        \begin{center}
\begin{tikzpicture}[scale=1.2]
    % 1. Dibujar el círculo
    \draw[thick, gray!70] (0,0) circle (2);

    % 2. Centro
    \filldraw[black] (0,0) circle (1.5pt);

    % 3. Coordenadas de los puntos
    \coordinate (P1) at (0:2);
    \coordinate (P2) at (100:2);
    \coordinate (P3) at (220:2);

    % 4. Triángulo
    \draw[thick, PosterHeading] (P1) -- (P2) -- (P3) -- cycle;

    % 5. Puntos coloreados y etiquetas
    \filldraw[blue!70!black]   (P1) circle (2.5pt) node[right=3pt] {$\phi_1$};
    \filldraw[orange!85!black] (P2) circle (2.5pt) node[above=3pt] {$\phi_2$};
    \filldraw[green!60!black]  (P3) circle (2.5pt) node[below left=3pt] {$\phi_3$};
\end{tikzpicture}
\end{center}
        
        \item Sea $\mathcal{C}$ un ciclo de $k$ aristas sin antípodas, luego:
        \[\#(E(\mathcal{C}) \cap\mathcal{T})>4|q(\mathcal{{C)}}|-k\].
    \end{itemize}
    
\end{borderbox}

% ========================================================
% 3. CONTINUACIÓN DE LA COLUMNA DERECHA
% ========================================================

\begin{borderbox}{Búsqueda SDP}

Utilizamos el número de winding para guiar la búsqueda externa hacia perturbaciones del estado \(C_4\).

\begin{equation*}
    \varphi_a(h) = \frac{\pi}{2}j(a) + hs_a
\end{equation*}

Donde \(j(a)\in\{0,1,2,3\}, h > 0\).

\begin{center}
  \begin{minipage}[t]{0.31\linewidth}
    \centering
    \PosterImage[0.9\linewidth]{assets/C4.png}
    \captionof{figure}{Estado \(C_4\).}
    \label{fig:c4}
  \end{minipage}%
  \hfill
  \begin{minipage}[t]{0.31\linewidth}
    \centering
    \PosterImage[0.9\linewidth]{assets/s7.png}
    \captionof{figure}{Estado \(S_7\).}
    \label{fig:s7}
  \end{minipage}%
  \hfill
  \begin{minipage}[t]{0.31\linewidth}
    \centering
    \PosterImage[0.9\linewidth]{assets/c4-certified-template-transparent.png}
    \captionof{figure}{\(T\) mejor obtenida}
    \label{fig:c8}
  \end{minipage}
\end{center}

Una búsqueda externa propone \((x,\varphi)\) y conexiones permitidas. Para cada geometría, resolvemos un problema semidefinido que maximiza la conectividad límite, sujeto a:

\begin{itemize}
    \item Equilibrio a orden cero y uno.
    \item Cancelación de la rigidez de primer orden.
    \item Estabilidad semidefinida(positiva) a segundo orden a margen \(\gamma\).
\end{itemize}

Finalmente, trasladamos los mejores candidatos a perturbaciones finitas y verificamos numéricamente su equilibrio y estabilidad.

\begin{center}
    Resultado: \(\boxed{\mu_0 \geq 0.69290770636958636365}\)
\end{center}

Nuestro trabajo generaliza el análisis perturbativo de \cite{okeeffe20261116densekuramotonetworks} a una búsqueda automatizada sobre distintas geometrías de clústeres.

\end{borderbox}

\begin{filledbox}{Conclusiones}
    \begin{itemize}
        \item Aumentar \(\mu\) por arriba de \(0,66\) se vuelve un problema de distribuir óptimamente las aristas en tensión por todo el grafo.
        \item Los ciclos que aumentan la cantidad de aristas en tensión en un régimen de \(q\) bajo son los de 4 aristas; conjeturamos que estos son los mejores para acercarse a \(\mu_c\).
        \item Se pudo obtener via busqueda SDP un resultado muy similar al mejor valor actual.
    \end{itemize}
\end{filledbox}

\begin{unfilledbox}{Referencias}
    % 1. Remove the automatic "References" title that breaks the box
    \renewcommand{\bibsection}{}
    
    % 2. Make the font significantly smaller for the poster's footer
    \bodyfont
    
    % 3. Reduce the vertical spacing between each citation
    \setlength{\bibsep}{10pt}
    
    % Generate the bibliography
    \bibliography{bibliography}
\end{unfilledbox}

\end{multicols*}
\end{document}
